% !TEX root = ../main.tex

\section{Silnik Carnota - sprawność w punkcie maksymalnej mocy}
\begin{figure}[h]
	\centering
	\begin{tikzpicture}[x=0.5cm, y=0.5cm]
	\tikzmath{
		\tankwidth = 10;
		\tanksep = 6;
		\tankthickness = 1;
		\engwidth = 1;
		\engsep = 5;
		\radius = 2;
		\chansep = 1;
		\auxtanksep = 5;
	}
	\draw[color=red, thick] (-\tankwidth, \tanksep) -- (\tankwidth, \tanksep) node[anchor=south, pos=0.5]{\(T_h\)};
	\draw[color=red, thick] (0, \auxtanksep) -- (\tankwidth, \auxtanksep) node[anchor=south, pos=0.5]{\(T_{hw}\)};
	\draw[color=red, thick] (-\tankwidth, \tanksep) -- (-\tankwidth, \tanksep+\tankthickness);
	\draw[color=red, thick] (\tankwidth, \tanksep) -- (\tankwidth, \tanksep+\tankthickness);
	\fill[color=red, opacity=0.2] (-\tankwidth, \tanksep) rectangle (\tankwidth, \tanksep+\tankthickness);
	
	\draw[color=blue, thick] (-\tankwidth, -\tanksep) -- (\tankwidth, -\tanksep) node[anchor=north, pos=0.5]{\(T_c\)};
	\draw[color=blue, thick] (0, -\auxtanksep) -- (\tankwidth, -\auxtanksep) node[anchor=north, pos=0.5]{\(T_{cw}\)};
	\draw[color=blue, thick] (-\tankwidth, -\tanksep) -- (-\tankwidth, -\tanksep-\tankthickness);
	\draw[color=blue, thick] (\tankwidth, -\tanksep) -- (\tankwidth, -\tanksep-\tankthickness);
	\fill[color=blue, opacity=0.2] (-\tankwidth, -\tanksep) rectangle (\tankwidth, -\tanksep-\tankthickness);
	\node[anchor=center] at (-\engsep, 0) {\(\eta_C\)};
	
	\draw[black] (-\engsep-\engwidth, -\tanksep) -- (-\engsep-\engwidth, \tanksep);
	\draw[black] (-\engsep+\engwidth, -\tanksep) -- (-\engsep+\engwidth, -\chansep-\radius);
	\draw[black] (-\engsep+\engwidth, \chansep+\radius) -- (-\engsep+\engwidth, \tanksep);
	\draw[black] (-\engsep+\engwidth, -\chansep-\radius) arc (180:90:\radius);
	\draw[black] (-\engsep+\engwidth, \chansep+\radius) arc (180:270:\radius);
	\node[anchor=center] at (\engsep, 0) {\(\eta_C'\)};
	
	\draw[black] (\engsep-\engwidth, -\auxtanksep) -- (\engsep-\engwidth, \auxtanksep);
	\draw[black] (\engsep+\engwidth, -\auxtanksep) -- (\engsep+\engwidth, -\chansep-\radius);
	\draw[black] (\engsep+\engwidth, \chansep+\radius) -- (\engsep+\engwidth, \auxtanksep);
	\draw[black] (\engsep+\engwidth, -\chansep-\radius) arc (180:90:\radius);
	\draw[black] (\engsep+\engwidth, \chansep+\radius) arc (180:270:\radius);
\end{tikzpicture}
	\caption{Schemat porównawczy silnika Carnota odwracalnego (po lewej) oraz endoodwracalnego (po prawej).}
\end{figure}
Teoretyczny cykl Carnota jest nieosiągalny. Gdy cykl ten przebiega z maksymalną sprawnością nie występuje w nim produkcja entropii wynikająca z przepływu ciepła. Stąd musi on zachodzić bez spadku temperatury między zbiornikiem ciepła a gazem w silniku. Taki przepływ jest nieskończenie wolny, więc moc takiego silnika jest zerowa \cite{curzon}.

Bardziej zbliżona do warunków rzeczywistych jest analiza silnika w punkcie, w którym moc jest maksymalna. Aby uwzględnić spadek temperatury na ściankach silnika rozpatrujemy warunki endoodwracalne, co oznacza, że czynnik roboczy jest w równowadze termodynamicznej wewnętrznie (gr. \'{e}ndon - wewnątrz), ale nie jest w równowadze z otoczeniem. W ogólności oznacza to, że występuje różnica temperatur między czynnikiem roboczym a zbiornikiem, z którym jest ten silnik w kontakcie.

Ilość ciepła przenikającego przez ścianki silnika jest proporcjonalne do różnicy temperatury:
\begin{gather}
\begin{split}
	Q_{1-2} &= \alpha_{c}\tau_{c}\pab{T_{c} - T_{cw}}, \\
	Q_{3-4} &= \alpha_{h}\tau_{h}\pab{T_{h} - T_{hw}},
\end{split}
\end{gather}
gdzie \(\alpha_{h}\) oraz \(\alpha_{c}\) to współczynniki przewodzenia ciepła zależne od materiału, grubości i powierzchni; a \(\tau_{h}\) oraz \(\tau_{c}\) to okresy czasu, podczas których następuje przepływ ciepła w jednym cyklu z odpowiednio zbiornika ciepłego i zimnego. Wprowadzamy wielkości pomocnicze zdefiniowane jako:
\begin{align}
	x &= T_h - T_{hw}, \\
	y &= T_{cw} - T_c.
\end{align}
Wykorzystując je, ciepła wymienione przez silnik są równe:
\begin{gather}\label{eq:endo_carnot_Q}
	\begin{split}
		Q_{1-2} &=  - \alpha_{c}\tau_{c}y, \\
		Q_{3-4} &=    \alpha_{h}\tau_{h}x.
	\end{split}
\end{gather}

Cykl Carnota działa tu między temperaturami \(T_{hw}\) i \(T_{cw}\) zamiast \(T_{h}\) i \(T_{c}\) jak było to w podrozdziale \ref{ssec:cykl_carnota}. Wewnętrzna produkcja entropii jest zerowa. Na cykl składają się 4 procesy: 2 izentropowe i 2 izotermiczne:
\begin{equation}\label{eq:endo_carnot_deltaS}
	\Delta S_{1-2} + \Delta S_{3-4} = 0, \quad\text{stąd}\quad
	\frac{Q_{1-2}}{T_{cw}} = - \frac{Q_{3-4}}{T_{hw}}.
\end{equation}
Podstawiając \eqref{eq:endo_carnot_Q} do \eqref{eq:endo_carnot_deltaS} otrzymujemy:
\begin{equation}\label{eq:endo_carnot_tau}
	\frac{\tau_{h}}{\tau_{c}} = \frac{\alpha_{c} y \pab{T_{h} - x}}{\alpha_{h} x \pab{T_{c} + y}}.
\end{equation}

Moc silnika jako funkcja \(x\) oraz \(y\) wyrażona jest wzorem:
\begin{equation}
	P\pab{x,y} = \frac{Q_{1-2} + Q_{3-4}}{\zeta\pab{\tau_{h} + \tau_{c}}},
\end{equation}
gdzie \(\zeta\) jest parametrem wiążącym długość całego cyklu \(\tau\) z sumą \(\tau_{h} + \tau_{c}\) równaniem:
\begin{equation}
	\tau = \zeta\pab{\tau_h + \tau_c}.
\end{equation}
Po podstawieniu \eqref{eq:endo_carnot_Q} moc ma postać:
\begin{equation}
	P\pab{x,y} =
	\frac{\alpha_{h}\tau_{h}x - \alpha_{c}\tau_{c}y}{\zeta\pab{\tau_h + \tau_c}} =
	\frac{\alpha_{h}\frac{\tau_h}{\tau_c}x - \alpha_{c}y}{\zeta\pab{\frac{\tau_h}{\tau_c} + 1}}.
\end{equation}
Po uwzględnieniu \eqref{eq:endo_carnot_tau} mamy \cite{curzon}:
\begin{equation}\label{eq:endo_carnot_P}
	P\pab{x,y} = \frac{\alpha_{h}\alpha_{c}xy}{\zeta} \frac{T_{h} - T_{c} - x - y}{\alpha_{h}T_{c}x + \alpha_{c}T_{h}y + \pab{\alpha_{h} - \alpha_{c}}xy}.
\end{equation}
Pierwsze pochodne cząstkowe mocy względem parametrów \(x\) oraz \(y\) są równe:
\begin{equation}\label{eq:endo_carnot_pdvPxy}
\begin{split}
	\pdv{P\pab{x,y}}{x} &= 
	\frac{\alpha_{h}\alpha_{c}}{\zeta} \frac{
	\alpha_{c}T_{h}y^2 \pab{T_{h}-T_{c}-x-y} - xy\bab{\alpha_{h}T_{c}x + \alpha_{c}T_{h}y + \pab{\alpha_{h}-\alpha_{c}}xy}
	}{\bab{\alpha_{h}T_{c}x + \alpha_{c}T_{h}y + \pab{\alpha_{h}-\alpha_{c}}xy}^2}, \\
	\pdv{P\pab{x,y}}{y} &= 
	\frac{\alpha_{h}\alpha_{c}}{\zeta} \frac{
	\alpha_{h}T_{c}x^2 \pab{T_{h}-T_{c}-x-y} - xy\bab{\alpha_{h}T_{c}x + \alpha_{c}T_{h}y + \pab{\alpha_{h}-\alpha_{c}}xy}
	}{\bab{\alpha_{h}T_{c}x + \alpha_{c}T_{h}y + \pab{\alpha_{h}-\alpha_{c}}xy}^2}.
\end{split}
\end{equation}

W punkcie maksymalnej mocy zerują się pierwsze pochodne cząstkowe względem \(x\) oraz \(y\). Otrzymujemy stąd warunki na wielkości w punkcie maksymalnej mocy \(x_0\) i \(y_0\):
\begin{align}\label{eq:endo_carnot_pdv0x}
	\alpha_{c}T_{h}\pab{T_{h} - T_{c} - x_0 - y_0}y_0 = x_0\bab{\alpha_{h}T_{c}x_0 + \alpha_{c}T_{h}y_0 + \pab{\alpha_{h} - \alpha_{c}}x_0y_0},
\end{align}
oraz:
\begin{align}\label{eq:endo_carnot_pdv0y}
	\alpha_{h}T_{c}\pab{T_{h} - T_{c} - x_0 - y_0}x_0 = y_0\bab{\alpha_{h}T_{c}x_0 + \alpha_{c}T_{h}y_0 + \pab{\alpha_{h} - \alpha_{c}}x_0y_0}.
\end{align}
Dzieląc stronami otrzymujemy:
\begin{align}\label{eq:endo_carnot_y}
	y_0^2 &= x_0^2\frac{\alpha_{h}T_{c}}{\alpha_{c}T_{h}}; & 
	y_0 &= x_0\sqrt{\frac{\alpha_{h}T_{c}}{\alpha_{c}T_{h}}}.
\end{align}
Podstawiając je do \eqref{eq:endo_carnot_pdv0x} oraz \eqref{eq:endo_carnot_pdv0y} i przekształcając otrzymujemy równanie kwadratowe w zmiennej \(\frac{x}{T_{h}}\):
\begin{equation}
	\pab{1 - \frac{\alpha_{h}}{\alpha_{c}}}\pab{\frac{x_0}{T_{h}}}^2 - 2\bab{\sqrt{\frac{\alpha_{h}T_{c}}{\alpha_{c}T_{h}}} + 1}\pab{\frac{x_0}{T_{h}}} + \pab{1 - \frac{T_{c}}{T_{h}}} = 0,
\end{equation}
które ma dwa rozwiązania rzeczywiste, ale tylko jedno z nich jest fizyczne i spełniające \(x_0/T_{h} < 1\). Następnie wykorzystując \eqref{eq:endo_carnot_y} mamy \cite{curzon}:
\begin{align}\label{eq:endo_carnot_xy}
	\frac{x_0}{T_{h}} &= \frac{1 - \sqrt{\frac{T_{c}}{T_{h}}}}{1 + \sqrt{\frac{\alpha_{h}}{\alpha_{c}}}}; & 
	\frac{y_0}{T_{c}} &= \frac{\sqrt{\frac{T_{h}}{T_{c}}} - 1}{1 + \sqrt{\frac{\alpha_{c}}{\alpha_{h}}}}.
\end{align}

Sprawność teoretycznego silnika Carnota działającego między zbiornikami o temperaturach \(T_h\) i \(T_c\) wynosi \(\eta_C = 1 - \frac{T_{c}}{T_{h}}\), podobnie dla rozpatrywanego działającego między temperaturami \(T_{hw}\) i \(T_{cw}\):
\begin{equation}\label{eq:endo_carnot_eta_def}
	\eta_C' = 1 - \frac{T_{cw}}{T_{hw}} = 
	1 - \frac{T_{c} + y_0}{T_{h} - x_0} = 
	1 - \frac{T_{c}}{T_{h}}\frac{1 + \frac{y_0}{T_{c}}}{1 - \frac{x_0}{T_{h}}}.
\end{equation}
Po podstawieniu \eqref{eq:endo_carnot_xy} i dłuższej manipulacji mamy:
\begin{equation}\label{eq:endo_carnot_eta}
	\eta_C' = 1 - \sqrt{\frac{T_{c}}{T_{h}}} = \frac{1 - \frac{T_{c}}{T_{h}}}{1 + \sqrt{\frac{T_{c}}{T_{h}}}} < 
	1 - \frac{T_{c}}{T_{h}} = \eta_C.
\end{equation}
Sprawność ta jest mniejsza od sprawności silnika idealnego, ale wciąż jest jedynie funkcją stosunku temperatur \(\frac{T_{c}}{T_{h}}\) i niezależna od innych parametrów opisujących silnik. Zerowanie się pierwszych pochodnych cząstkowych nie jest warunkiem wystarczającym istnienia maksimum. Aby roztrzygnąć czy jest to maksimum, minimum czy punkt siodłowy przeprowadzono dodatkowe obliczenia w dodatku \ref{sec:endo_carnot_max}.

Z porównania \eqref{eq:endo_carnot_eta_def} oraz \eqref{eq:endo_carnot_eta} otrzymujemy równość:
\begin{equation}\label{eq:endo_carnot_frac_Tw}
	\frac{T_{cw}}{T_{hw}} = \sqrt{\frac{T_{c}}{T_{h}}},
\end{equation}
a z podstawienia \eqref{eq:endo_carnot_xy} do \eqref{eq:endo_carnot_tau} oraz wykorzystując \eqref{eq:endo_carnot_frac_Tw} mamy:
\begin{equation}
	\frac{\tau_{h}}{\tau_{c}} = \sqrt{\frac{\alpha_{c}}{\alpha_{h}}}.
\end{equation}
Przekształcając równanie \eqref{eq:endo_carnot_xy} można znaleźć:
\begin{align}\label{eq:endo_carnot_Tw}
	T_{hw} &= \sqrt{T_{h}} \frac{\sqrt{\alpha_{h}T_{h}} + \sqrt{\alpha_{c}T_{c}}}{\sqrt{\alpha_{h}} + \sqrt{\alpha_{c}}}; & 
	T_{cw} &= \sqrt{T_{c}} \frac{\sqrt{\alpha_{h}T_{h}} + \sqrt{\alpha_{c}T_{c}}}{\sqrt{\alpha_{h}} + \sqrt{\alpha_{c}}}.
\end{align}
Wkorzystując równania \eqref{eq:endo_carnot_Tw} maksymalna moc wyrażona jest wzorem:
\begin{equation}\label{eq:endo_carnot_Pmax}
	P = \frac{\alpha_{h}\tau_{h}\pab{T_{h} - T_{hw}} - \alpha_{c}\tau_{c}\pab{T_{cw} - T_{c}}}{\zeta\pab{\tau_{h} + \tau_{c}}} = \frac{\alpha_{h}\alpha_{c}}{\zeta}\pab{\frac{\sqrt{T_{h}} - \sqrt{T_{c}}}{\sqrt{\alpha_{h}} + \sqrt{\alpha_{c}}}}^2.
\end{equation}
Wyprowadzony wzór na moc w punkcie maksymalnej mocy jest zależny jedynie od temperatur zbiorników, współczynników opisujących przepływ ciepła oraz współczynnika czasowego \(\zeta\).
